% Generated draft: requires mathematical and source review.
\documentclass[11pt]{article}
\usepackage[margin=25mm]{geometry}
\usepackage{fontspec}
\setmainfont{Latin Modern Roman}
\usepackage{amsmath,amssymb,mathtools,unicode-math}
\setmathfont{Latin Modern Math}
\usepackage{enumitem,needspace,xurl,hyperref}
\setlength{\parindent}{0pt}
\setlength{\parskip}{5pt}
\setlength{\emergencystretch}{3em}
\newcommand{\N}{\mathbb{N}}
\newcommand{\Z}{\mathbb{Z}}
\newcommand{\Q}{\mathbb{Q}}
\newcommand{\R}{\mathbb{R}}
\newcommand{\C}{\mathbb{C}}
\newcommand{\F}{\mathbb{F}}
\newcommand{\A}{\mathbb{A}}
\newcommand{\PP}{\mathbb{P}}
\newcommand{\E}{\mathbb{E}}
\newcommand{\eps}{\varepsilon}
\newcommand{\dd}{\,\mathrm{d}}
\begin{document}
\section*{UnsolvedMath ID: 2659}
\textbf{Source problem code:} catalogue entry
\subsection{Definitions and mathematical statement}
An elementary abelian group is a vector space over a prime field \(\F_p\), regarded additively; its rank is its vector-space dimension, finite or infinite. For a group \(H\), let \(Z(H)=\{h\in H:hu=uh\text{ for every }u\in H\}\) and let \(H'=[H,H]\) be its derived subgroup.

The source considers an extension \(1\to A\to G\to B\to1\), with \(A\) and \(B\) elementary abelian, and an element \(a\in A\) with trivial stabilizer \(C_B(a)\). Since \(A\) is abelian, conjugation on \(A\) factors through \(G/A\cong B\), so that stabilizer is meaningful. To form the displayed subgroup \(\langle a,B\rangle\) literally, however, a copy of \(B\) inside \(G\) is needed. In the split case \(G=A\rtimes B\), as in the supplied counterexample, use its embedded complement and put \(H=\langle a,B\rangle\). The proposed inequality is
\[
C_B(a)=\{1\}\quad\Longrightarrow\quad
\operatorname{rank}(Z(H)\cap H')\leq\operatorname{rank}(B).
\]
Here \(H'\leq A\), so the left group is elementary abelian and its rank is defined. No restriction to finite ranks is imposed by the short statement. An arbitrary nonsplit-extension reading would require the meaning of \(\langle a,B\rangle\) or a choice of lifts to be specified; the split case already falls under the intended inequality.

The supplied background records Monticone's negative answer in a letter of 21 March 2026. Let
\[
R=\F_3[x,y]/(x,y)^3,\qquad A=(R,+),
\]
and let \(B\leq\operatorname{Aut}(A)\) be generated by multiplication by \(1+x\) and by \(1+y\). Form \(G=A\rtimes B\), and take \(a=1\in R\), the additive-group element represented by the constant polynomial. The source reports
\[
B\cong C_3^2,\qquad C_B(a)=\{1_B\},\qquad
Z(\langle a,B\rangle)\cap\langle a,B\rangle'
=\langle x^2,xy,y^2\rangle_{\F_3}\cong C_3^3.
\]
Thus its reported ranks are 2 and 3, respectively. The stated quotient uses the cube of the ideal generated by \(x,y\), not the ideal generated only by \(x^3,y^3\).

Both registry and supplied literature triage dated 2026-08-17 classify the inequality solved negatively. The triage says the explicit example matches all intended hypotheses and reports an accompanying Lean formalization, a later joint paper and the original solution note. This draft reports those source claims without executing the formalization or independently checking the rank computation; proposed corrected bounds remain a separate question.
\subsection{Short English statement}
Does an element with trivial stabilizer in an elementary-abelian extension force the central derived subgroup of the generated group to have rank no larger than the quotient? The supplied characteristic-three split extension is reported to have quotient rank 2 and central derived rank 3, refuting the inequality.
\subsection{Sources}
\begin{itemize}
\item Local UnsolvedMath registry, permanent ID 2659, source code catalogue entry, ``Kourovka Notebook Problem 21.150,'' from \emph{Kourovka Notebook}, New Problems, 21st issue (2026), original problem 21.150; attributed to V. I. Zenkov. Includes supplied literature triage dated 2026-08-17.
\item Pietro Monticone, ``Solution of Kourovka Notebook Problem 21.150'' (2026), associated in the background with a letter of 21 March 2026: \url{https://kourovkanotebookorg.wordpress.com/wp-content/uploads/2026/03/solution_21_150-3.pdf}.
\item Wouter van Doorn, Elias Judin, Pietro Monticone, and Daniel Morrison, ``On Some Problems from the Kourovka Notebook,'' arXiv:2607.17477 (2026): \url{https://arxiv.org/abs/2607.17477}.
\item Pietro Monticone et al., Kourovka formalization repository, Problem 21.150 (2026): \url{https://github.com/pitmonticone/Kourovka}.
\end{itemize}
\end{document}
